\documentclass[11pt,a4paper]{article}

\usepackage[margin=2.5cm]{geometry}
\usepackage{amsmath}
\usepackage{booktabs}
\usepackage{xcolor}
\usepackage{hyperref}
\usepackage{float}
\usepackage{microtype}

\hypersetup{colorlinks=true, linkcolor=blue!70!black}

\title{\textbf{Design and Optimization of an Accelerator\\
for a CNN}\\[0.4em]
\large HW-SW Codesign --- Midterm Report}
\author{Hugo Svolgaard}
\date{Due: 2026-04-02\null}

\begin{document}

\maketitle

%% ============================================================
\section*{Network Overview}
%% ============================================================

The application runs a 5-layer CNN (256$\times$256 RGB input) through 5 sequential
Conv2D $\to$ MaxPool pairs followed by 2 dense layers.
Input data is stored in CHW layout (Channel $\times$ Height $\times$ Width).
Each conv layer's input channel count equals the previous layer's filter count;
the first layer starts from 3 (RGB).

\begin{center}
\begin{tabular}{ccccc}
\toprule
Layer & $C_i$ & $C_o$ & Input & Output \\
\midrule
0 & 3   & 32  & 256$\times$256 & 254$\times$254 \\
1 & 32  & 64  & 127$\times$127 & 125$\times$125 \\
2 & 64  & 128 & 62$\times$62   & 60$\times$60   \\
3 & 128 & 256 & 30$\times$30   & 28$\times$28   \\
4 & 256 & 64  & 14$\times$14   & 12$\times$12   \\
\bottomrule
\end{tabular}
\end{center}

%% ============================================================
\section{Step 1: Basic Accelerator --- Bias and ReLU in Hardware}
%% ============================================================

The HLS accelerator \texttt{Conv2D\_HW} mirrors the software 6-nested-loop structure.
Ports are exposed via \texttt{m\_axi} (Master AXI, bulk DDR transfers) and
\texttt{s\_axilite} (Slave AXI-Lite, scalar parameters and base addresses).
Pointer arguments appear on both: \texttt{s\_axilite} provides the runtime base
address; \texttt{m\_axi} with \texttt{offset=slave} performs the actual transfer from
that address. \texttt{input}/\texttt{output} and \texttt{filters}/\texttt{biases} are
split across two AXI master bundles (HP0 and HP1) so the inner loop can read a filter
coefficient and an input pixel simultaneously, achieving II=1.

\noindent In the software version, bias addition and ReLU are separate post-processing passes
over the full output tensor, each requiring an additional DDR read/write. In the
hardware version, both are folded directly into the accumulation step, eliminating
those extra passes entirely. ReLU is parameterised so the same accelerator handles
all conv layers (ReLU=1) and the final dense layer (ReLU=0).

%% ============================================================
\section{Step 2: Caching of Filter Coefficients}
%% ============================================================

In the minimal implementation, all $C_i \times 9$ filter coefficients are re-fetched
from DDR on every MAC operation (\mbox{outH $\times$ outW} times per filter).
The fix: copy the full filter into a local BRAM array (\texttt{localFilters[256*9]})
once before the spatial loops, using a pipelined sequential loop (II=1) that HLS
can infer as a single AXI burst. All subsequent MAC reads come from BRAM (1-cycle latency).
DDR filter reads drop from $N_f \times H \times W \times C_i \times 9$
to $N_f \times C_i \times 9$. The array is reused (overwritten) for each filter;
2304 elements covers the largest case (layer 4: $256 \times 9$).

\noindent AXI burst parameters were also tuned: \texttt{max\_read\_burst\_length=256} raises
throughput from 64 to 1024 bytes per burst; \texttt{num\_read\_outstanding=16} allows
16 requests in flight simultaneously, hiding DDR latency.

%% ============================================================
\section{Step 3: Caching of Input Rows}
%% ============================================================

The 3$\times$3 kernel requires 3 consecutive input rows per output row.
When the y-loop advances by 1, two of those rows overlap with the previous iteration
--- only 1 new row needs loading from DDR\@.
Three row buffers (\texttt{rowBuffer[3][4096]}, each holding $C_i \times W_{in}$
values) form a circular buffer. After the initial 3-row load, each y-iteration fetches
exactly 1 new row into the expired slot (\texttt{newRow \% 3}).

\noindent \texttt{\#pragma HLS LOOP\_FLATTEN off} is applied to the row-loading nested loops.
Without it, HLS merges the channel and pixel loops into a single flattened pipeline;
the resulting combined address computation (large inter-channel strides in CHW layout)
lengthens the critical path beyond the 10\,ns clock period, causing negative timing slack.
Keeping the loops separate costs a few cycles per channel boundary but preserves timing
closure.

%% ============================================================
\section{Step 4: Parallelising Output Filters}
%% ============================================================

For any given input pixel, computing it against $N$ output filters simultaneously
is fully independent: each filter reads the same pixel but multiplies by different
coefficients. The solution replicates the MAC hardware $N$ times:

\begin{itemize}
  \item \texttt{localFilters[N][256*9]}: N separate BRAM arrays;
        \texttt{ARRAY\_PARTITION complete dim=1} gives each a physically independent
        block so all N can be read in the same cycle.
  \item The inner \texttt{p} loop is \texttt{UNROLL}ed, instantiating N multipliers.
  \item Filter loading is restructured to interleaved (outer index \texttt{i},
        inner unrolled \texttt{p}): all N BRAMs are written simultaneously from
        a single sequential DDR read stream, eliminating N separate burst restarts.
  \item \texttt{nActive} guards the last group when \texttt{numFilters \% N $\neq$ 0}.
\end{itemize}

\noindent With N=16 and the inner loop reading one pixel and 16 filter values per cycle,
output writes to 16 scattered DDR planes became the bottleneck
(Conv 0 had only 27 MACs per pixel but 16 write transactions).
An \textbf{output row buffer} (\texttt{outRowBuf[16][256]}) accumulates one full row
per filter in BRAM, then burst-writes each row to DDR, converting 16 scattered
writes per pixel into 16 sequential bursts per output row. This alone gave a
7.2$\times$ speedup on Conv 0.

\noindent The \texttt{latency=20} AXI hint (which caused pessimistic pipeline bubbles) was
also removed, allowing HLS to schedule transactions more aggressively.

%% ============================================================
\section{Step 5 (Optional): Row Reuse and Overlapping Transfers}
%% ============================================================

\subsection*{5a: 4-Buffer Prefetch}

With 3 row buffers, the slot for the next prefetch (\texttt{(y+3)\%3}) collides
with the current computation slot (\texttt{y\%3}), forcing them to serialise.
A 4th buffer resolves the collision: \texttt{(y+3)\%4} is always different from
\texttt{y\%4}, \texttt{(y+1)\%4}, \texttt{(y+2)\%4}.
An additional benefit: \texttt{\%4} is a 2-bit mask in hardware vs.\ a full
division circuit for \texttt{\%3}, reducing the compute pipeline iteration latency
from 43 to 11 cycles (II=1 maintained).

\subsection*{5b: 2$\times$ X-position Parallelism}

Adjacent output x-positions share the \emph{same} filter coefficient
(filter address depends on \texttt{(ch,cy,cx)}, not \texttt{x}).
The x-loop steps by 2 with dual accumulators (\texttt{acc0}, \texttt{acc1}),
reading two adjacent input pixels and one filter value per cycle --- doubling MAC
throughput without extra filter BRAM.
\texttt{rowBuffer} and \texttt{outRowBuf} are partitioned \texttt{cyclic factor=2 dim=2}
so adjacent addresses land in different BRAM banks, enabling conflict-free parallel reads/writes.
DSP utilisation rose from 41\% to 63\%.

\subsection*{Negative slack fix}

Loop-flattening (removing \texttt{LOOP\_FLATTEN off}) was tried to reduce
per-channel burst restart overhead but caused negative timing slack.
\texttt{LOOP\_FLATTEN off} was re-applied selectively to the row-loading and
output-write loops while keeping the interleaved filter-loading structure.

%% ============================================================
\section{Bonus: MaxPool Hardware Accelerator}
%% ============================================================

After all Conv2D optimisations, Conv and MaxPool were equally bottlenecked (43.5\%
and 43.1\% of runtime respectively). MaxPool 2$\times$2 with stride 2 requires only
comparisons (zero DSPs), making it ideal for a lightweight accelerator.

\noindent The HLS design reads two consecutive rows in a single burst (they are contiguous in
CHW layout), applies a comparison tree at II=1, and burst-writes the output row.
Two additional HP ports (HP2, HP3) handle MaxPool input/output independently from the
Conv2D ports. The software proxy \texttt{CMaxPoolProxy} reuses Conv2D's DMA address
translation table via \texttt{ShareDMAMappings} () to avoid a double-free.

%% ============================================================
\section{Results}
%% ============================================================

Resource cost:
\[
  C_R = \frac{25}{140}\cdot\#\text{BRAMs}
      + \frac{25}{220}\cdot\#\text{DSPs}
      + \frac{25}{53200}\cdot\#\text{LUTs}
      + \frac{25}{106400}\cdot\#\text{FFs}
\]

\begin{table}[H]
\centering
\small
\setlength{\tabcolsep}{5pt}
\begin{tabular}{p{5.8cm}rrrrrrrr}
\toprule
\textbf{Description} & \textbf{Time (ms)} & \textbf{MHz} & \textbf{LUTs} & \textbf{FFs} & \textbf{BRAMs} & \textbf{DSPs} & \textbf{$C_R$} & \textbf{Pareto?} \\
\midrule
SW only (ARM)                       & 28000 & n/a & n/a   & n/a   & n/a & n/a & 0     & Yes \\
Initial HW                          & 86357 & 100 & 5316  & 5838  & 0   & 39  & 9.11  & No  \\
+ Dual AXI + Filter cache           & 14124 & 100 & 9180  & 5977  & 8   & 30  & 11.26 & No  \\
+ Input row caching                 & 12870 & 100 & 10369 & 6494  & 40  & 15  & 12.87 & No  \\
+ N=4 filter parallelism            & 4853  & 100 & 15986 & 12978 & 64  & 66  & 22.97 & No  \\
+ N=8 filter parallelism            & 3374  & 100 & 19224 & 15686 & 96  & 88  & 30.58 & No  \\
+ 4-buffer prefetch                 & 3190  & 100 & 17117 & 12607 & 96  & 88  & 28.74 & No  \\
+ N=16 + output row buffering       & 1402  & 100 & 17099 & 13367 & 176 & 91  & 42.09 & No  \\
+ Loop flatten + interleaved load   & 1397  & 100 & 16976 & 20575 & 176 & 92  & 43.22 & No  \\
+ 2$\times$ x-position parallelism  & 1116  & 100 & 24920 & 20913 & 192 & 140 & 51.18 & No  \\
+ MaxPool HW accelerator (initial)  & 744   & 100 & n/a   & n/a   & n/a & n/a & n/a   & n/a \\
\midrule
\textbf{Final: Conv2D (N=16, 2x x-par.) + MaxPool HW (merged burst)} & \textbf{739} & 100 & n/a & n/a & n/a & n/a & n/a & n/a \\
\bottomrule
\end{tabular}
\caption{Execution time and resource usage per solution (256$\times$256 input, layer 0: $C_i$=3, $C_o$=32).
% TODO: Pareto optimal analysis
% A solution is Pareto-optimal if no other solution achieves both strictly lower time AND
% strictly lower cost C_R. Fill in the Pareto? column based on Time vs C_R tradeoff.
% Plot Time (ms) vs C_R to identify the Pareto frontier.
}
\end{table}

\noindent The final solution accelerates all 5 Conv2D and all 5 MaxPool layers on the FPGA.
The remaining bottleneck is the 2 dense layers running in software on the ARM CPU (20\% of runtime).

\begin{table}[H]
\centering
\begin{tabular}{lrr}
\toprule
\textbf{Component} & \textbf{Time (ms)} & \textbf{\% of total} \\
\midrule
Conv2D (5 layers, HW)  & 512  & 69.3\% \\
MaxPool (5 layers, HW) & 77   & 10.4\% \\
Dense (2 layers, SW)   & 149  & 20.2\% \\
Flatten + Sigmoid      & 1    & 0.1\%  \\
\midrule
\textbf{Total}         & \textbf{739} & 100\% \\
\bottomrule
\end{tabular}
\caption{Final solution runtime breakdown. Overall speedup vs.\ software-only: \textbf{37.9$\times$} (28.0\,s to 0.739\,s).}
\end{table}

\end{document}
